\documentclass[11pt,a4paper]{article}

\usepackage{amsmath,amsthm,amssymb}
\usepackage{booktabs}
\usepackage[ruled,vlined,linesnumbered]{algorithm2e}
\usepackage[colorlinks=true,linkcolor=blue,citecolor=blue,urlcolor=blue]{hyperref}
\usepackage{geometry}
\geometry{margin=2.5cm}
\usepackage{microtype}

% ---------- theorem environments ----------
\newtheorem{theorem}{Theorem}[section]
\newtheorem{proposition}[theorem]{Proposition}
\newtheorem{corollary}[theorem]{Corollary}
\newtheorem{definition}[theorem]{Definition}
\newtheorem{remark}[theorem]{Remark}
\newtheorem{example}[theorem]{Example}

% ---------- notation macros ----------
\newcommand{\R}{\mathbb{R}}
\newcommand{\bq}{\mathbf{q}}
\newcommand{\bp}{\mathbf{p}}
\newcommand{\bu}{\mathbf{u}}
\newcommand{\bw}{\mathbf{w}}
\newcommand{\bg}{\mathbf{g}}
\newcommand{\bt}{\mathbf{t}}
\newcommand{\bz}{\mathbf{z}}
\newcommand{\bs}{\mathbf{s}}
\newcommand{\bk}{\mathbf{k}}
\newcommand{\bC}{\mathbf{C}}
\newcommand{\bS}{\mathbf{S}}
\newcommand{\bK}{\mathbf{K}}
\newcommand{\bM}{\mathbf{M}}
\newcommand{\bI}{\mathbf{I}}
\newcommand{\bF}{\mathbf{F}}
\newcommand{\bdelta}{\boldsymbol{\delta}}
\newcommand{\btau}{\boldsymbol{\tau}}
\newcommand{\bzero}{\mathbf{0}}
\newcommand{\khat}{\hat{\bk}}
\newcommand{\qhat}{\hat{\bq}}
\newcommand{\norm}[1]{\|#1\|}
\newcommand{\abs}[1]{\lvert#1\rvert}
\newcommand{\tr}{\top}
\newcommand{\Estar}{E^{*}}
\newcommand{\code}[1]{\texttt{#1}}

% ---------- algorithm2e style ----------
\SetAlgoLined
\DontPrintSemicolon
\SetKwComment{Comment}{// }{}

% =========================================================
\begin{document}

\title{\textbf{Frictional Contact on the Elastic Half-Space}\\[0.5em]
       \large Tangential Kernels, the Incremental Stick--Slip Problem,\\
       Projected Conjugate Gradients, and the Bipotential Reference Solver}
\author{ASPHER project (formerly Hcontact) --- technical reference}
\date{July 2026}
\maketitle

\begin{abstract}
\noindent
This note documents the frictional (tangential) contact capability of the
ASPHER boundary-element contact solver. It is written to be readable by
someone new to computational contact mechanics: every object is introduced
from its physical meaning, small worked examples with numbers accompany the
key ideas, and the algorithms are stated exactly as implemented (file names
are given in the margin of the text as \code{typewriter} references). The
frictionless normal problem and its Polonsky--Keer projected conjugate
gradient (CG) solver are documented in the companion note \code{pcg.tex};
the fast operators (H$^2$/FMM, FFT convolution) in
\code{h2\_fmm\_detailed.tex}. Here we cover: the Cerruti tangential
elasticity kernels and their element integrals
(\S\ref{sec:cerruti}), friction laws and the threshold field
(\S\ref{sec:laws}), the quasi-static incremental stick--slip problem and its
formulation as a convex quadratic program (\S\ref{sec:incremental}), the
production solver --- a vector projected CG with a two-metric projection,
spectral preconditioning, and an outer Newton iteration for force control
(\S\ref{sec:solver}), the incremental driver with history and dissipation
tracking (\S\ref{sec:driver}), the de~Saxc\'e--Feng bipotential Uzawa
reference solver (\S\ref{sec:uzawa}), and the analytical benchmarks
(Cattaneo--Mindlin, Mindlin stiffness, Ciavarella--J\"ager) used as
validation gates (\S\ref{sec:benchmarks}).
\end{abstract}

\tableofcontents

% ---------------------------------------------------------
\section{Introduction: What Friction Adds to Contact}
\label{sec:intro}

\subsection{A tabletop experiment}

Press your hand flat on a table and push it sideways, gently at first. Two
things happen in sequence. While the push is small, the hand does not move:
the interface transmits a \emph{shear traction} (a tangential force per unit
area) that exactly balances your push. This is \emph{stick}. Increase the
push and at some point the hand starts to slide: the interface cannot
transmit more shear than a threshold, and beyond it the surfaces move
relative to each other. This is \emph{slip}. The classical
\emph{Coulomb law} says the threshold is proportional to how hard the
surfaces are pressed together: slip starts when the shear traction $q$
reaches $\mu p$, where $p$ is the local contact pressure and $\mu$ is the
friction coefficient (typically $0.1$--$1$ for dry engineering surfaces).

The interesting --- and mathematically rich --- regime is the one in
between, called \emph{partial slip}: because elastic bodies deform, some
parts of a contact patch can slip while others still stick. Press a rubber
eraser on paper and twist or push it slightly: the edge of the contact
slips (you may see wear marks) while the centre remains stuck. Partial slip
is not a curiosity; it governs fretting fatigue, the damping of bolted
joints, the stiffness of wheel--rail and tyre--road contacts, and the
nucleation of earthquakes on faults.

Two consequences of the threshold behaviour drive the entire structure of
this document:

\begin{enumerate}
\item \textbf{History dependence.} Where slip has occurred, the interface
  ``remembers'' it: shear tractions do not return to zero when the load is
  removed (they become \emph{locked in}). The state of the interface depends
  on the whole loading \emph{path}, not just the current load. Hence we must
  solve a \emph{sequence of load increments}, carrying state from step to
  step (\S\ref{sec:driver}).
\item \textbf{Dissipation.} Slip against friction converts mechanical work
  into heat. Every load increment must dissipate a non-negative energy
  $D \ge 0$; this is both a physical statement (second law) and a sharp
  correctness test for the solver (\S\ref{sec:dissipation}).
\end{enumerate}

\subsection{From the frictionless to the frictional problem}

The frictionless solver (see \code{pcg.tex}) finds the contact pressure
$p(x,y) \ge 0$ between a rigid rough surface and an elastic half-space,
using the Boussinesq relation ``pressure $\to$ normal surface displacement''
and a projected CG. The frictional extension adds the tangential direction:

\begin{center}
\begin{tabular}{lll}
\toprule
 & normal (existing) & tangential (new) \\
\midrule
traction unknown       & $p$ (scalar, $\ge 0$)  & $\bq = (q_x, q_y)$ (2-vector, $\abs{\bq} \le s$) \\
surface displacement   & $u_z = \bS p$            & $\bu_t = (u_x,u_y) = \bC \bq$ \\
rigid-body unknown     & approach $\alpha$        & tangential shift $\bdelta_t = (\delta_x,\delta_y)$ \\
elastic kernel         & Boussinesq/Love          & Cerruti (\S\ref{sec:cerruti}) \\
inequality law         & no adhesion, no overlap  & stick--slip friction (\S\ref{sec:laws}) \\
constraint on unknown  & $p_i \ge 0$              & $\abs{\bq_i} \le s_i$ (disk) \\
solver                 & scalar projected CG      & vector projected CG (\S\ref{sec:solver}) \\
\bottomrule
\end{tabular}
\end{center}

The structure is deliberately parallel: the tangential problem is again a
convex quadratic program (QP) with simple per-point constraints, solved by
the same family of projected CG methods, accelerated by the same kind of
FFT spectral preconditioner, on the same fast operators (FFT convolution or
H$^2$/FMM).

\subsection{Scope and assumptions (v1)}
\label{sec:scope}

\begin{itemize}
\item \textbf{Uncoupled elasticity.} For contact between \emph{similar}
  materials (equal elastic constants, or the special combinations with
  Dundurs' $\beta = 0$), pressure produces no tangential surface
  displacement and shear produces no normal displacement. The normal and
  tangential elastic problems then decouple: $u_z = \bS p$ and
  $\bu_t = \bC \bq$ with no cross terms. ASPHER v1 assumes this; the
  coupled (dissimilar-material) case is a documented follow-up, for which
  the bipotential solver of \S\ref{sec:uzawa} is the designated strategy.
  The friction \emph{law} still couples the two problems one way: the
  slip threshold depends on $p$.
\item \textbf{Quasi-static, small displacements.} Loads change slowly
  (no inertia); the contact geometry is not advected by the tangential
  shift (the gap field is fixed within a step). Time enters only through
  the slip \emph{velocity} in rate-dependent friction laws, via the step
  duration $\Delta t$.
\item \textbf{Rigid rough counterpart.} As in the normal problem, one body
  is rigid and rough, the other is an elastic half-space with effective
  modulus $\Estar$ and Poisson ratio $\nu$; the roughness enters through
  the initial gap field $g_0(x,y)$.
\end{itemize}

\subsection{Notation}

Discretisation: an $N_s \times N_s$ uniform grid of square elements of side
$h = L/N_s$, $N = N_s^2$ points. Per-point 2-vectors are written in bold
with a subscript $i$: $\bq_i = (q_{x,i}, q_{y,i})$. The full tangential
traction vector is stored \emph{stacked},
$\bq = [\,q_x\,;\,q_y\,] \in \R^{2N}$ (all $x$-components first). The
tangential influence operator $\bC \in \R^{2N \times 2N}$ is symmetric
positive definite (SPD) with $2 \times 2$ block structure
$\bigl(\begin{smallmatrix} \bC_{xx} & \bC_{xy} \\ \bC_{xy} & \bC_{yy}
\end{smallmatrix}\bigr)$. $\qhat_i = \bq_i / \abs{\bq_i}$ is the unit vector
along $\bq_i$. Superscript $n$ marks the previous accepted load step;
$\Delta(\cdot)$ is the increment over the current step. $G$ is the shear
modulus, related to the effective modulus by $\Estar (1-\nu) = 2G$ for the
rigid--elastic pair (used repeatedly below).

% ---------------------------------------------------------
\section{Tangential Elasticity: the Cerruti Kernels}
\label{sec:cerruti}

\subsection{Reminder: the normal (Boussinesq--Love) problem}

The half-space is the standard idealisation for contact: when the contact
patch is small compared to the bodies, each body responds as if it were an
elastic solid filling a half of space, loaded on its flat boundary. For a
point force $P$ pressing \emph{normally} on the surface at the origin,
Boussinesq's solution gives the normal displacement of the surface at
distance $\rho$:
\begin{equation}
  u_z(\rho) = \frac{P}{\pi \Estar \rho}.
\end{equation}
Superposition over a distributed pressure $p(x',y')$ turns this into a
convolution, $u_z = \bS p$; integrating the $1/\rho$ kernel exactly over
each square element gives Love's closed form and the discrete influence
matrix used by ASPHER (see \code{pcg.tex} \S6 and
\code{boussinesq\_kernel.cpp}).

\subsection{Cerruti's point-force solution}

Cerruti~\cite{Cerruti1882} solved the companion problem: a point force $Q_x$
applied \emph{tangentially} to the surface (along $x$). On the surface
($z=0$), at position $(x,y)$ with $\rho = \sqrt{x^2 + y^2}$, the tangential
surface displacements are
\begin{equation}
  u_x = \frac{Q_x}{2\pi G}\left[ \frac{1-\nu}{\rho}
        + \nu\,\frac{x^2}{\rho^3} \right],
  \qquad
  u_y = \frac{Q_x}{2\pi G}\; \nu\,\frac{x y}{\rho^3}.
  \label{eq:cerruti-point}
\end{equation}
Three features deserve attention before any algebra:
\begin{itemize}
\item Like Boussinesq's, the kernel decays as $1/\rho$ --- the tangential
  problem has exactly the same long-range character, so all the fast-operator
  machinery (H-matrices, FMM, FFT) applies unchanged.
\item Unlike the normal case, the response is \emph{anisotropic}: the
  displacement along the force ($u_x$) differs from the transverse response,
  and a force along $x$ produces a (smaller, $\propto \nu$) displacement
  along $y$. The influence operator is therefore a $2\times2$ tensor.
\item At $\nu = 0$ the coupling term vanishes and
  $u_x = Q_x/(2\pi G \rho)$: the tangential problem becomes two independent
  copies of the normal problem up to a constant factor. This limit is the
  key to the sharpest validation tests (\S\ref{sec:cj}).
\end{itemize}
The prefactor identity used throughout the code is worth one line: with
$\Estar = E/(1-\nu^2)$ (rigid counterpart) and $2G = E/(1+\nu)$,
\begin{equation}
  \Estar(1-\nu) = \frac{E(1-\nu)}{1-\nu^2} = \frac{E}{1+\nu} = 2G
  \qquad\Longrightarrow\qquad
  \frac{1}{2\pi G} = \frac{1}{\pi \Estar (1-\nu)} .
  \label{eq:prefactor}
\end{equation}

\subsection{Element-integrated kernels (Pohrt--Li closed forms)}
\label{sec:element}

As in the normal problem, we do not use the singular point kernel directly:
each element carries a \emph{uniform} traction, and the influence
coefficient is the exact integral of \eqref{eq:cerruti-point} over the
source rectangle $[-a,a]\times[-b,b]$. Closed forms are taken from
Pohrt \& Li~\cite{PohrtLi2014} (with Dydo \& Busby~\cite{DydoBusby1995} as
an independent source). Introduce the corner shorthands
\begin{equation}
  k = \abs{x}+a,\quad l = \abs{x}-a,\quad m = \abs{y}+b,\quad n = \abs{y}-b,
  \qquad R(u,v) = \sqrt{u^2+v^2},
\end{equation}
and the two ``log brackets''
\begin{align}
  X_{\log} &= k \ln\!\frac{m + R(k,m)}{n + R(k,n)}
            + l \ln\!\frac{n + R(l,n)}{m + R(l,m)}, \\
  Y_{\log} &= m \ln\!\frac{k + R(k,m)}{l + R(l,m)}
            + n \ln\!\frac{l + R(l,n)}{k + R(k,n)}.
\end{align}
Then the three rectangle integrals of the point-kernel ingredients are
\begin{equation}
  \int\!\!\int \frac{dA'}{\rho} = X_{\log} + Y_{\log}
  \quad\text{(this is Love's normal bracket)},\qquad
  \int\!\!\int \frac{(x-x')^2}{\rho^3}\, dA' = Y_{\log},
\end{equation}
\begin{equation}
  \int\!\!\int \frac{(x-x')(y-y')}{\rho^3}\, dA'
  = R(k,n) - R(k,m) + R(l,m) - R(l,n).
\end{equation}
Combining with \eqref{eq:cerruti-point}, the element-integrated tangential
displacements due to unit $q_x$ on the rectangle are
\begin{equation}
  \boxed{\;
  u_x = \frac{(1-\nu)\,X_{\log} + Y_{\log}}{2\pi G},
  \qquad
  u_y = \frac{\nu}{2\pi G}
        \bigl[ R(k,n) - R(k,m) + R(l,m) - R(l,n) \bigr].
  \;}
  \label{eq:element}
\end{equation}
(The $u_x$ bracket is $(1-\nu)(X_{\log}+Y_{\log}) + \nu\,Y_{\log}
= (1-\nu)X_{\log} + Y_{\log}$.) These are \code{cerruti\_uxx} and
\code{cerruti\_uxy} in \code{cerruti\_kernel.cpp}. Loading along $y$ follows
by the $x \leftrightarrow y$ symmetry.

\begin{remark}[A typo in the literature, and why we verify]
Equation (18) of \cite{PohrtLi2014} as printed carries spurious $h^2$
factors, dimensionally inconsistent with their own eq.~(13). The corner form
in \eqref{eq:element} is the correct rectangle integral of
$x'y'/\rho^3$; it was verified against adaptive $64^2$-point Gauss--Legendre
quadrature of the point kernel (relative error $< 10^{-9}$ on all sampled
offsets, including self-terms; test \code{test\_cerruti}) and against an
independent thesis derivation (Duquesne 2025, Annexe~G, in the project's
reference collection \code{REF/Tangential/}), whose corner primitives reduce
to \eqref{eq:element} exactly. The lesson generalises: \emph{never} transcribe
closed-form kernels from a paper without an independent numerical
cross-check.
\end{remark}

\begin{example}[The self-term, and the Mindlin compliance ratio]
\label{ex:selfterm}
Evaluate \eqref{eq:element} at the centre of the loaded square element
($x=y=0$, $a=b=h/2$). Then $k=m=a$, $l=n=-a$, all four $R$'s equal
$a\sqrt2$, and both log brackets collapse to
$X_{\log} = Y_{\log} = 4a \ln(1+\sqrt2)$, giving
\begin{equation}
  u_x^{\text{self}} = \frac{(2-\nu)\, 2h \ln(1+\sqrt2)}{2\pi G},
  \qquad u_y^{\text{self}} = 0 .
\end{equation}
Compare the normal self-term $u_z^{\text{self}} = 4h\ln(1+\sqrt2)/(\pi\Estar)$:
using \eqref{eq:prefactor},
\begin{equation}
  \frac{u_x^{\text{self}}}{u_z^{\text{self}}}
  = \frac{2-\nu}{2(1-\nu)}
  \;\approx\; 1.21 \quad (\nu = 0.3).
\end{equation}
The surface is about 20\% \emph{more} compliant tangentially than normally
--- precisely Mindlin's classical ratio of tangential to normal contact
compliance~\cite{Mindlin1949,Johnson1985}. A one-element sanity check thus
already reproduces a landmark result of contact mechanics.
\end{example}

\subsection{Symmetries, storage, and the discrete operator}
\label{sec:symmetries}

On the uniform grid the influence coefficients depend only on the offset
$(\Delta_x, \Delta_y)$ between target and source elements (translation
invariance), so each of the three components is an $N_s\times N_s$ lookup
table, as for Love. Parities cut storage further
(\code{CerrutiKernel}):
\begin{itemize}
\item $K_{xx}(\Delta_x,\Delta_y)$ is even in both arguments;
\item $K_{yy}(\Delta_x,\Delta_y) = K_{xx}(\Delta_y,\Delta_x)$ (the
  $x\leftrightarrow y$ transpose --- not stored);
\item $K_{xy}$ is odd in $\Delta_x$ and odd in $\Delta_y$ (stored for
  non-negative offsets; signs restored at lookup; exactly zero on the axes,
  which the lookup guards explicitly).
\end{itemize}
The dense operator (used only for testing) is the $2N\times2N$ SPD matrix
with blocks $[\,\bC_{xx}, \bC_{xy}; \bC_{xy}, \bC_{yy}\,]$ acting on the
stacked $[\,q_x; q_y\,]$ layout.

\subsection{The Fourier symbol}
\label{sec:symbol}

For periodic or whole-plane problems, convolution operators are diagonal in
Fourier space; the ``symbol'' (the Fourier multiplier) is the single most
useful analytical object we have --- it drives the preconditioner
(\S\ref{sec:precond}), validates the tables, and estimates operator norms
(\S\ref{sec:uzawa}). The 2-D surface Fourier transforms of the point-kernel
ingredients are
\begin{equation}
  \widehat{\left(\tfrac1\rho\right)} = \frac{2\pi}{\abs\bk},\qquad
  \widehat{\left(\tfrac{x^2}{\rho^3}\right)} = \frac{2\pi k_y^2}{\abs\bk^3},\qquad
  \widehat{\left(\tfrac{xy}{\rho^3}\right)} = -\frac{2\pi k_x k_y}{\abs\bk^3},
\end{equation}
which assemble into the $2\times2$ matrix symbol of $\bC$
(\code{CerrutiKernel::symbol}):
\begin{equation}
  \boxed{\;
  \widehat\bC(\bk) = \frac{2}{\Estar (1-\nu)\,\abs\bk}
  \left[ \bI - \nu\, \khat\,\khat^{\tr} \right],
  \qquad \khat = \bk/\abs\bk .
  \;}
  \label{eq:symbol}
\end{equation}
Its eigenvectors are the directions along and transverse to the wave vector:
\begin{equation}
  \widehat\bC\,\khat = \underbrace{\frac{2}{\Estar \abs\bk}}_{\text{longitudinal}} \khat,
  \qquad
  \widehat\bC\,\khat^{\perp} = \underbrace{\frac{2}{\Estar (1-\nu)\abs\bk}}_{\text{transverse}} \khat^{\perp}.
\end{equation}
The longitudinal eigenvalue equals \emph{exactly} the Love symbol
$2/(\Estar\abs\bk)$ of the normal operator --- a strong internal consistency
check tying the tangential constants to the long-validated normal path, and
the identity that makes the $\nu = 0$ limit collapse onto the normal
problem. The discrete tables reproduce \eqref{eq:symbol} to the same
truncation accuracy as the Love table reproduces its symbol
(\code{test\_cerruti}: diagonal modes to $0.6$--$1.2\%$, dominated by the
finite-table truncation, not by the closed forms).

\subsection{Fast operators: FFT convolution and H$^2$/FMM}
\label{sec:operators}

Both production operators from the normal problem carry over
(\code{tangential\_operator.cpp}); the dense matrix exists for tests only.

\paragraph{FFT convolution (\code{TangentialFFTOperator}).}
The same zero-padded (Hockney) scheme as the scalar \code{FFTOperator}:
the three kernel tables are embedded in wrap-around order on the padded
$(2N_s)^2$ grid and transformed once at build. A useful subtlety: although
$K_{xy}$ is odd in each axis separately, all three tables are even under
\emph{joint} negation $(\Delta_x,\Delta_y) \to (-\Delta_x,-\Delta_y)$, so
all three spectra are \emph{real} (matching the real continuum symbol
\eqref{eq:symbol}; asserted at build). One matvec is then: two zero-pad
scatters $\to$ two forward FFTs $\to$ an all-real per-mode $2\times2$ mix
\begin{equation}
  \widehat u_x = \widehat K_{xx} \widehat q_x + \widehat K_{xy} \widehat q_y,
  \qquad
  \widehat u_y = \widehat K_{xy} \widehat q_x + \widehat K_{yy} \widehat q_y,
\end{equation}
$\to$ two inverse FFTs $\to$ two gathers. This operator is \emph{exact}
(matches the dense matvec to $2.9\times10^{-16}$ relative $L_2$;
\code{test\_tangential}) and is what the production driver uses.

\paragraph{H$^2$/FMM (\code{TangentialH2Operator}).}
Three scalar \code{H2Operator} instances (xx, yy, xy) built through the
kernel-functor constructor --- the far field interpolates the smooth
off-origin Cerruti kernels with Chebyshev polynomials exactly as it does
Love. A blocked matvec costs four scalar applies
($u_x = \bC_{xx} q_x + \bC_{xy} q_y$, etc.). Accuracy mirrors the normal
path: relative $L_2$ vs dense $1.7\times10^{-4}$ at Chebyshev order $q{=}4$,
$7.0\times10^{-6}$ at $q{=}6$. Preferred at very large $N_s$ (O($N$) memory).

% ---------------------------------------------------------
\section{Friction Laws and the Threshold Field}
\label{sec:laws}

\subsection{The threshold abstraction}

All friction laws in ASPHER reduce to one object: a per-point scalar
\emph{threshold field} $s_i \ge 0$, the largest shear traction magnitude
point $i$ can sustain,
\begin{equation}
  s_i = \tau_c\!\left(p_i,\; \abs{\dot\bw_i},\; T_i\right),
\end{equation}
which may depend on the local pressure $p_i$, the slip velocity magnitude
$\abs{\dot\bw_i}$, and a temperature (or any auxiliary) field $T_i$. The
implemented models (\code{friction\_model.hpp}) are:
\begin{center}
\begin{tabular}{lll}
\toprule
model & threshold & typical use \\
\midrule
Tresca & $s_i = \tau_c$ where $p_i > 0$, else $0$ &
  analysis, limits, plasticity-like interfaces \\
Coulomb & $s_i = \mu \max(p_i, 0)$ &
  the standard dry-friction law \\
generic callback & $s_i = \tau_c(p_i, v_i, T_i)$, user function &
  rate-and-state-like, thermal weakening, \dots \\
\bottomrule
\end{tabular}
\end{center}
Two contractual properties are enforced no matter what the user law
returns: $s_i \ge 0$ everywhere, and $s_i = 0$ wherever $p_i \le 0$ (no
shear transmission without contact --- an open point cannot carry
traction). The callback wrapper sanitises its output accordingly (clamps
negatives and NaNs to zero, zeroes non-contact points, size-checks), so no
user law can hand the solver an inadmissible threshold field.

\subsection{Cone and disk}
\label{sec:cone}

Geometrically, the Coulomb law confines the full traction vector
$(q_x, q_y, p)$ at each point to the convex \emph{friction cone}
\begin{equation}
  K_\mu = \left\{ (\bq_t, p) : \abs{\bq_t} \le \mu\, p \right\},
\end{equation}
a cone in 3-D traction space with its axis along $p$ and half-opening
$\arctan\mu$. The production solver (\S\ref{sec:solver}) is
\emph{staggered}: it solves the normal problem first, freezes $p$, and then
solves the tangential problem. With $p_i$ (hence $s_i$) frozen, the cone
constraint at each point becomes a horizontal slice --- a \emph{disk}
$\abs{\bq_i} \le s_i$ in the $(q_x,q_y)$ plane. Disks are the tangential
analogue of the normal problem's half-line $p_i \ge 0$, and ``project onto
the disk'' (radially rescale a too-long vector back to length $s_i$)
replaces ``clamp negative pressures to zero''. The bipotential solver
(\S\ref{sec:uzawa}) instead works with the full cone directly and never
staggers.

Within one increment, Coulomb with frozen $p$ \emph{is} a Tresca problem
with spatially varying threshold. The genuinely Coulomb-specific coupling
(slip threshold changing as $p$ changes) re-enters across load steps
through the driver.

% ---------------------------------------------------------
\section{The Quasi-Static Incremental Problem}
\label{sec:incremental}

\subsection{Kinematics and sign conventions}
\label{sec:kinematics}

Let $\bdelta_t$ be the rigid tangential shift of the (rigid) counterpart and
$\bu_{t,i} = (\bC\bq)_i$ the elastic tangential displacement of surface
point $i$. The \emph{local slip} of the counterpart relative to the surface
point is
\begin{equation}
  \bw_i = \bdelta_t - (\bC \bq)_i .
  \label{eq:slip}
\end{equation}
Sign convention (the single most error-prone spot in any friction code; it
is fixed here once and used consistently, including in the bipotential sign
map of \S\ref{sec:uzawa}): $\bq_i$ is the traction exerted \emph{on the
half-space}. Friction \emph{drags the surface along with the slipping
counterpart}, so during slip the traction is parallel to the slip,
$\bq_i \parallel +\bw_i$ (and the reaction on the counterpart opposes its
motion, as intuition demands).

The friction law, per point, on the total traction:
\begin{equation}
  \underbrace{\abs{\bq_i} < s_i \;\Rightarrow\; \dot\bw_i = \bzero}_{\text{stick}},
  \qquad
  \underbrace{\abs{\bq_i} = s_i \;\text{and}\; \dot\bw_i = \lambda_i \qhat_i,\ \lambda_i \ge 0}_{\text{slip}},
  \qquad
  \underbrace{p_i = 0 \;\Rightarrow\; s_i = 0 \;\Rightarrow\; \bq_i = \bzero}_{\text{open}}.
  \label{eq:law}
\end{equation}
Quasi-statics has no natural time scale, so ``rates'' are discretised by
backward Euler over a load step $n \to n{+}1$ of duration $\Delta t$: the
\emph{slip increment} is
\begin{equation}
  \Delta\bw_i = \Delta\bdelta_t - \bigl(\bC\,(\bq - \bq^n)\bigr)_i,
  \qquad \text{slip velocity } \abs{\Delta\bw_i}/\Delta t .
  \label{eq:slipinc}
\end{equation}
$\Delta t$ matters only for velocity-dependent laws; for Tresca/Coulomb the
steps are pure sequence, no clock.

\subsection{Warm-up: a one-point toy problem}
\label{sec:toy1}

Strip the problem to one point and one direction: a single ``elastic
asperity'' of stiffness $k$ (traction per unit displacement, standing in for
$\bC^{-1}$), threshold $s$, dragged by a base shift $\delta$. The traction
is $q = k(\delta - w)$ where $w$ is accumulated slip. Start from rest and
increase~$\delta$:
\begin{equation}
  q(\delta) = \min(k\delta,\, s), \qquad
  w(\delta) = \max\!\left(0,\ \delta - s/k\right).
\end{equation}
The point sticks until $\delta = s/k$, then slips with constant traction.
Now reverse the shift after loading to $\delta_{\max} > s/k$. The point
first re-sticks (elastic unloading, $q = s - k(\delta_{\max}-\delta)$) and
only re-slips backwards when $q$ reaches $-s$, i.e.\ after a reversal of
$2s/k$ --- \emph{twice} the forward stick range. Three lessons scale up
unchanged to the full field problem:
\begin{enumerate}
\item \textbf{Hysteresis}: the $(\delta, q)$ curve on unloading is not the
  loading curve reversed; the enclosed loop area is dissipated energy.
\item \textbf{Locked-in traction}: choose $k = 1$, $s = 1$, load to
  $\delta = 2$ ($q = 1$, $w = 1$), then unload to $\delta = 1$: the point
  sticks the whole way back and ends with $q = 0$ but $w = 1$ ---
  \emph{zero traction, non-zero permanent slip}. Unloading to $\delta = 0$
  ends with $q = -1$: a locked-in \emph{negative} traction at zero applied
  shift. The interface state cannot be reconstructed from the current load.
\item \textbf{Path dependence has structure}: for \emph{monotonic}
  proportional loading the response is history-free (this is what makes the
  Ciavarella--J\"ager superposition of \S\ref{sec:cj} possible); history
  effects appear on unloading and reloading.
\end{enumerate}

\subsection{The incremental quadratic program}
\label{sec:qp}

For a \emph{frozen} threshold field $s$ (one Tresca-like pass; the
velocity-dependent outer loop is \S\ref{sec:threshold-loop}), the load
step's tangential problem is the convex QP
\begin{equation}
  \boxed{\;
  \min_{\bq}\ J(\bq) = \tfrac12\,(\bq - \bq^n)^\tr \bC\, (\bq - \bq^n)
                       \;-\; \Delta\bdelta_t \cdot \sum_i \bq_i
  \qquad \text{s.t.}\quad \abs{\bq_i} \le s_i \ \ \forall i .
  \;}
  \label{eq:qp}
\end{equation}
Under \textbf{displacement control} $\Delta\bdelta_t$ is prescribed. Under
\textbf{force control} $\Delta\bdelta_t$ is unknown and the mean traction is
constrained instead, $\frac1N \sum_i \bq_i = \bar\bq$ (equivalently, the
total tangential force $Q = h^2 \sum_i \bq_i = L^2 \bar\bq$); $\bdelta_t$
then plays the role of the Lagrange multiplier of the load constraint,
exactly as the rigid approach $\alpha$ does in the normal problem.

Where does \eqref{eq:qp} come from? It is the standard incremental
variational statement for this class of dissipative laws: the first
(elastic) term is the energy of the traction \emph{increment}, and the
constraint set encodes the maximal admissible traction. One does not need
that pedigree to use it, though, because its optimality conditions can be
checked directly against the friction law --- which we now do, since it is
also exactly how the solver's convergence metric is built.

\subsection{KKT conditions $\Leftrightarrow$ stick--slip law}
\label{sec:kkt}

The gradient of $J$ at point $i$ is
\begin{equation}
  \bg_i = \nabla_{\bq_i} J = \bigl(\bC(\bq - \bq^n)\bigr)_i - \Delta\bdelta_t
        = -\,\Delta\bw_i
  \label{eq:gradient}
\end{equation}
by \eqref{eq:slipinc}: \emph{the residual of the QP is minus the slip
increment}. Form the Lagrangian with multipliers $\lambda_i \ge 0$ for the
disk constraints $\abs{\bq_i} - s_i \le 0$:
\begin{equation}
  \mathcal L = J(\bq) + \sum_i \lambda_i \left( \abs{\bq_i} - s_i \right).
\end{equation}
Stationarity at a point with $\bq_i \ne \bzero$ gives
$\bg_i + \lambda_i \qhat_i = \bzero$. Reading off the two cases:
\begin{itemize}
\item \textbf{Interior} ($\abs{\bq_i} < s_i$): complementary slackness
  forces $\lambda_i = 0$, so $\bg_i = \bzero$, i.e.\
  $\Delta\bw_i = \bzero$ --- the point \emph{sticks}. \checkmark
\item \textbf{On the bound} ($\abs{\bq_i} = s_i$):
  $\Delta\bw_i = -\bg_i = \lambda_i \qhat_i$ with $\lambda_i \ge 0$ --- the
  point \emph{slips}, by a non-negative amount, \emph{parallel to} $+\qhat_i$.
  \checkmark
\end{itemize}
So the KKT conditions of \eqref{eq:qp} are precisely the incremental form of
the friction law \eqref{eq:law}: solving the QP \emph{is} enforcing
friction. Uniqueness of $\bq$ follows from strict convexity ($\bC$ is SPD).
This equivalence also dictates the solver's error metric
(\S\ref{sec:metric}): at an interior point any non-zero $\bg_i$ is a
violation; at a bound point the violations are the component of $\bg_i$
\emph{perpendicular} to $\qhat_i$ (slip not aligned with traction) and any
\emph{positive} component along $\qhat_i$ (a wrong-sign multiplier ---
slip opposing the traction).

\subsection{Dissipation}
\label{sec:dissipation}

The frictional work of one increment (per unit area $h^2$ of each element)
is
\begin{equation}
  D = h^2 \sum_i \bq_i \cdot \Delta\bw_i .
\end{equation}
At the QP solution: stick points contribute $0$ ($\Delta\bw_i = \bzero$),
slip points contribute
$\bq_i \cdot \lambda_i \qhat_i = \lambda_i \abs{\bq_i} = \lambda_i s_i \ge 0$.
Hence
\begin{equation}
  D = h^2 \sum_{i \,\in\, \text{slip}} \lambda_i\, s_i \;\ge\; 0 ,
\end{equation}
with equality iff no slip occurs. The driver computes and reports $D$ every
step, and the test suite asserts $D \ge 0$ (up to the solver's floor) on
every scenario --- a cheap, global, physics-grounded regression net that has
no analogue in the frictionless problem.

\subsection{A two-point worked example}
\label{sec:toy2}

The smallest problem that exhibits \emph{partial slip} has two points and
elastic coupling. Take one direction only ($q_y \equiv 0$), compliance
matrix and thresholds
\begin{equation}
  \bC = \begin{pmatrix} 2 & 1 \\ 1 & 2 \end{pmatrix},
  \qquad s = (1,\ 3),
\end{equation}
(in arbitrary consistent units), no history ($\bq^n = 0$), displacement
control with $\delta = 3.6$.

\emph{Try full stick}: $\bC \bq = \delta \mathbf 1$ gives
$\bq = \delta\,\bC^{-1}\mathbf 1 = (1.2,\ 1.2)$. Point 2 is fine
($1.2 \le 3$) but point 1 violates its bound ($1.2 > 1$): full stick is
inconsistent, the contact \emph{must} partially slip.

\emph{Correct partition}: point 1 slips at its bound $q_1 = s_1 = 1$;
point 2 sticks, so its slip is zero:
$w_2 = \delta - (q_1 + 2 q_2) = 0 \Rightarrow q_2 = (3.6 - 1)/2 = 1.3 \le 3$
\checkmark. Verify point 1's slip direction:
$w_1 = \delta - (2 q_1 + q_2) = 3.6 - 3.3 = 0.3 > 0$, same sign as $q_1$
\checkmark. Both KKT cases check out; dissipation
$D = q_1 w_1 = 0.3 > 0$.

Note what the elastic coupling did: the slipping point 1 sheds load onto its
sticking neighbour ($q_2 = 1.3 >$ the full-stick $1.2$). In a real contact
this load shedding is what makes the stick zone shrink smoothly from the
edges inward as the load grows --- the Cattaneo--Mindlin annulus of
\S\ref{sec:cm} is exactly this mechanism with $N$ large.

\subsection{Velocity-dependent thresholds: the outer fixed point}
\label{sec:threshold-loop}

If $\tau_c$ depends on the slip velocity, the threshold field depends on the
solution: $s = \tau_c(p, \abs{\Delta\bw(\bq)}/\Delta t, T)$. The driver
resolves this with a damped fixed point: solve the QP with the current
$s^{(k)}$, evaluate the law at the resulting slip rate, and average,
\begin{equation}
  s^{(k+1)} = \tfrac12\left( s^{(k)}
            + \tau_c\!\left(p, \abs{\Delta\bw^{(k)}}/\Delta t, T\right)\right),
\end{equation}
warm-starting each inner QP from the previous pass, until the relative
change of $s$ drops below \code{threshold\_rtol} (default $10^{-3}$).
Rate-independent laws take exactly one pass; a rate-weakening smoke test
converges in 7 damped passes. (The velocity in pass 0 is taken from the
\emph{previous step's} slip increment --- quasi-static continuation.)

% ---------------------------------------------------------
\section{The Production Solver: Vector Projected CG}
\label{sec:solver}

\subsection{The staggered step}

Because the elasticity is uncoupled (\S\ref{sec:scope}), each load step
splits cleanly (\code{friction\_driver.cpp}):
\begin{enumerate}
\item \textbf{Normal solve}: the existing Polonsky--Keer PCG at the target
  mean pressure, warm-started from the previous step. $p$ never feels $\bq$.
\item \textbf{Threshold}: $s = \tau_c(p, \cdot, T)$, with the outer loop of
  \S\ref{sec:threshold-loop} if velocity-dependent.
\item \textbf{Tangential QP} \eqref{eq:qp}: the vector projected CG below.
\item \textbf{State update}: commit $\bq$, $\bu_t$, $\bdelta_t$, accumulated
  slip, dissipation (transactionally --- \S\ref{sec:transactional}).
\end{enumerate}

The tangential solver (\code{friction\_solve.cpp}, entry
\code{solve\_tangential}) is the vector generalisation of Polonsky--Keer:
same skeleton (projected CG + exact line search + restart logic), with the
scalar bound $p_i \ge 0$ replaced by the disk $\abs{\bq_i} \le s_i$. Three
design points differ from a naive transcription, each forced by a measured
failure mode; they are worth understanding because they generalise to any
projected CG on disk-like constraints.

\subsection{Candidate set and disk projection}

The CG iterates over the full \emph{candidate set}
$\mathcal A = \{ i : s_i > 0 \}$ (all points that can carry any traction),
\emph{not} over the current stick set. After each CG step the update is
projected point-wise onto the disks (radial clamp):
\begin{equation}
  \bq_i \leftarrow
  \begin{cases}
    \bq_i, & \abs{\bq_i} \le s_i,\\[2pt]
    s_i\, \bq_i / \abs{\bq_i}, & \abs{\bq_i} > s_i .
  \end{cases}
\end{equation}
Why not iterate only on the stick set, mimicking the normal solver's
free-set logic? Because a bound point is not ``done'': it still owns a
\emph{direction} degree of freedom. $\abs{\bq_i} = s_i$ pins the magnitude,
but $\qhat_i$ must rotate until it aligns with the slip --- and a stick-set
CG never updates it. Keeping all candidate points in the CG lets slip
directions converge; the clamp enforces the magnitudes. (This matches the
approach of Tamaas' \code{PolonskyKeerTan}.)

\subsection{The two-metric projection}
\label{sec:twometric}

At a bound (slipping) point the residual is, by \S\ref{sec:kkt},
\begin{equation}
  \bg_i = -\lambda_i\, \qhat_i + \bg_i^{\perp},
\end{equation}
an outward-radial\footnote{The CG convention here is
$\bq \leftarrow \bq - \tau\, \bt$ with $\bt$ built from $\bg$; a component
of $\bt$ along $-\qhat_i$ moves $\bq_i$ \emph{outward}, i.e.\ against the
constraint.} multiplier part plus a tangential misalignment part. The
multiplier part is not an error --- at the exact solution it is the only
thing left --- yet a naive CG feeds it into the search direction, where it
does pure harm: the outward push is annihilated by the clamp, but it has
already polluted the \emph{exact line search}, whose $\tau$ minimises the
objective along the unclamped ray. Measured effect: objective oscillation
and stall at $\sim 10^{-2}$ relative error. The remedy, classical under the
name \emph{two-metric projection}~\cite{Bertsekas1982}, is to strip the
offending component from the direction source $\bz$ at every bound point:
\begin{equation}
  \zeta_i = \bz_i \cdot \qhat_i < 0
  \quad\Longrightarrow\quad
  \bz_i \leftarrow \bz_i - \zeta_i\, \qhat_i ,
\end{equation}
keeping inward components ($\zeta_i > 0$: the point may return to stick)
and all transverse components (direction rotation). The projection
(clamping) and the direction metric thus differ --- hence the name.

Two companion safeguards, both from observed failures:
\begin{itemize}
\item \textbf{$\beta$-restart on partition change only.} Conjugacy is reset
  ($\beta = 0$) when a point changes stick/slip status --- the objective's
  active geometry changed. But \emph{re-clamping an already-bound point is
  not a partition change}: while slip directions rotate, points are
  re-clamped every iteration, and restarting $\beta$ there degrades CG to
  steepest descent (measured: stall at $\sim 10^{-2}$).
\item \textbf{Descent fallback.} If the preconditioned, conjugated direction
  loses descent ($\bg^\tr \bt \le 0$), one retry is made with the plain
  projected residual (steepest descent); without the retry the iteration can
  deadlock recomputing the same non-descent direction forever.
\end{itemize}

\subsection{Line search, error metric, and best-iterate return}
\label{sec:metric}

The step length is the exact minimiser on the candidate set,
$\tau = (\bg^\tr\bt)\,/\,(\bt^\tr \bC\, \bt)$, with all grid-length
reductions accumulated in double (the float-safety conventions of
\code{pcg.tex} \S6 apply verbatim). The convergence metric is the
complementarity form of the KKT conditions of \S\ref{sec:kkt}:
\begin{equation}
  e \;=\; \frac{1}{\bigl(\sum_i s_i\bigr)\, g_{\text{scale}}}
  \left[
  \sum_{i\,\text{interior}} (s_i - \abs{\bq_i})\, \abs{\bg_i}
  \;+\;
  \sum_{i\,\text{bound}} s_i \left( \abs{\bg_i^{\perp}}
      + \max(0,\ \bg_i\cdot\qhat_i) \right)
  \right],
  \label{eq:metric}
\end{equation}
the direct analogue of Polonsky--Keer's $\sum p_i \abs{g_i}$: every term
vanishes exactly at the solution and each measures a specific law violation
(interior slip; misaligned slip; wrong-sign multiplier).
$g_{\text{scale}}$ is the first-iteration residual scale (optionally floored
--- see \S\ref{sec:warm}).

Near its floor ($\sim 10^{-5}$--$10^{-4}$ relative on hard partial-slip
problems --- the clamps perturb CG conjugacy, and rotating slip directions
keep injecting small residual) the exact line search becomes
noise-dominated and the iterate can wander \emph{away} from the best point
visited (measured: error touches $7\times10^{-6}$, then oscillates at
$10^{-3}$--$10^{-2}$). The solver therefore tracks and returns the
\emph{best iterate seen}, and a stagnation guard (200 non-improving
iterations) converts a floor-limited run into an honest converged-at-floor
exit, mirroring \code{solve\_contact}.

\begin{algorithm}[t]
\caption{Vector projected CG, displacement control
  (\code{solve\_disp} in \code{friction\_solve.cpp})}
\label{alg:pcgtan}
\KwIn{operator $\bC$ (matvec), thresholds $s \ge 0$, shift $\Delta\bdelta_t$,
      optional warm start $\bq^{(0)}$, history $\bu_{\mathrm{hist}}$, tol}
$\mathcal A \leftarrow \{ i : s_i > 0\}$;\quad
clamp $\bq^{(0)}$ onto the disks; $\beta_{\mathrm{on}} \leftarrow 0$\;
\For{$k = 0, 1, 2, \dots$}{
  $\bu \leftarrow \bC \bq$;\quad
  $\bg_i \leftarrow \bu_i + \bu_{\mathrm{hist},i} - \Delta\bdelta_t$
  on $\mathcal A$ \Comment{$\bg = -\Delta\bw$}
  evaluate error \eqref{eq:metric}; track best iterate; stop on tol or
  200-iteration stagnation\;
  $\bz \leftarrow \bM^{-1}\bg$ on $\mathcal A$
  \Comment{spectral preconditioner \S\ref{sec:precond}, or $\bz = \bg$}
  \textbf{two-metric strip} at bound points:
  $\bz_i \mathrel{-}= \min(0, \bz_i\cdot\qhat_i)\, \qhat_i$\;
  $\beta \leftarrow \beta_{\mathrm{on}} \cdot
    \max\bigl(0, \bz^\tr(\bg - \bg_{\mathrm{prev}})\bigr) / G_{\mathrm{old}}$
    \Comment{PR+; FR optional}
  $\bt \leftarrow \bz + \beta\, \bt$;\quad
  $\tau \leftarrow (\bg^\tr\bt) / (\bt^\tr \bC \bt)$
  \Comment{fallback to steepest descent if $\bg^\tr\bt \le 0$}
  $\bq_i \leftarrow \Pi_{\mathrm{disk}}\!\left( \bq_i - \tau\, \bt_i \right)$
  on $\mathcal A$ \Comment{radial clamp}
  $\beta_{\mathrm{on}} \leftarrow 0$ if the stick/slip partition changed,
  else $1$\;
}
\Return best iterate, stick/slip/open classification, $\bq$-mean\;
\end{algorithm}

\subsection{The $2\times2$ spectral preconditioner}
\label{sec:precond}

The normal solver's iteration count is tamed by preconditioning with the
inverse symbol $\abs\bk$ (see \code{pcg.tex} \S6); the tangential operator
needs the matrix version. From \eqref{eq:symbol},
$\widehat\bC \propto \abs\bk^{-1} (\bI - \nu \khat\khat^\tr)$. Since
$\khat\khat^\tr$ is a projector ($P^2 = P$), the inverse of $\bI - \nu P$
is found by the ansatz $\bI + \gamma P$:
\begin{equation}
  (\bI - \nu P)(\bI + \gamma P) = \bI + (\gamma - \nu - \nu\gamma) P
  = \bI
  \iff
  \gamma = \frac{\nu}{1-\nu},
\end{equation}
so (dropping the overall scale, which cancels in CG)
\begin{equation}
  \boxed{\;
  \bM^{-1}(\bk) = \abs\bk \left( \bI + \frac{\nu}{1-\nu}\,
  \khat\khat^{\tr} \right).
  \;}
\end{equation}
Application (\code{TangentialFourierPreconditioner} in
\code{fourier\_precond.cpp}) follows the scalar preconditioner exactly:
mask the stacked residual to the candidate set, two forward FFTs, the
per-mode symmetric $2\times2$ multiply, two inverse FFTs, masked gather;
the per-component mean is removed only under force control (where the load
constraint, not CG, fixes the mean). Like its scalar sibling this is an
\emph{approximate} inverse (non-padded, masked), affecting only direction
quality --- never the converged solution. One implementation subtlety: at
the self-conjugate Nyquist modes the odd off-diagonal entry $w_{xy}$ must
be zeroed (an odd real symbol cannot be represented there). Measured on an
$N_s = 64$ partial-slip problem: $271 \to 126$ iterations ($2.15\times$),
identical solutions at the KKT level.

\subsection{Force control: outer Newton on the rigid shift}
\label{sec:force}

Under force control the constraint $\operatorname{mean}(\bq) = \bar\bq$
must be met with $\bdelta_t$ as its multiplier. The scalar normal solver
enforces its load constraint \emph{inside} the CG loop by an additive
correction each iteration; the direct transcription of that trick to
friction is \textbf{unstable} and was abandoned after measurement: the
additive load correction, the interior-mean estimate of $\bdelta_t$, and
the exact line search form a feedback loop with period-2 oscillation
(measured gain $\approx 1.5$ per iteration at 40\% slip --- each correction
overshoots the previous one). The stable architecture separates concerns:

\begin{enumerate}
\item \textbf{Inner}: displacement-controlled solves
  (Algorithm~\ref{alg:pcgtan}) at a given shift $\bdelta$.
\item \textbf{Outer}: solve the 2-dimensional root problem
  \begin{equation}
    \bF(\bdelta) = \operatorname{mean}\bigl(\bq(\bdelta)\bigr) - \bar\bq
    = \bzero
  \end{equation}
  by Newton's method with a Broyden-updated $2\times2$ stiffness $\bK$.
  The map $\bdelta \mapsto \operatorname{mean}(\bq)$ is monotone (SPD
  full-stick stiffness, softening as slip develops), so Newton from the
  \emph{full-stick} stiffness under-shoots and approaches the root stably.
\end{enumerate}

Supporting machinery, each piece answering a measured failure:
\begin{itemize}
\item \textbf{Probe initialisation.} $\bK$ is initialised from two
  small-shift displacement solves ($\varepsilon$-probes along $x$ and $y$;
  $F(\bzero) = -\bar\bq$ needs no solve since $\bdelta = \bzero \Rightarrow
  \bq = \bzero$ cold). The probe scale comes from one matvec on the
  gross-slip traction field; if the probes slip more than $20\%$,
  $\varepsilon$ is halved once. With history present (\S\ref{sec:warm}) the
  probes are taken \emph{differentially} around the operating point instead.
\item \textbf{Noise-guarded Broyden updates.} The inner solver resolves
  $\operatorname{mean}(\bq)$ only to its metric floor; a secant computed
  from sub-floor differences is pure noise and can collapse $\bK$
  (observed: $\det \bK$ decaying $10^{-1} \to 10^{-17} \to$ Newton
  blow-up). Updates are rejected unless the updated determinant stays
  safely non-degenerate, and the outer loop \textbf{stops when $\bF$ stops
  responding} to $\bdelta$ at the inner resolution (floor detection).
\item \textbf{Terminal exact-load correction.} The best-$\abs\bF$ iterate is
  kept, and the residual load ($\sim 10^{-4}\,\abs{\bar\bq}$, the inner
  floor) is then distributed \emph{additively over the interior (stick)
  points} --- the same correction that is unstable inside the loop is
  perfectly safe as a one-shot post-step (two or three re-clamping passes
  reach machine precision; final load met to $\sim 3\times10^{-16}$
  relative). The KKT residual absorbs a perturbation of the same harmless
  order.
\item \textbf{Gross-slip guard.} No equilibrium exists if
  $\abs{\bar\bq} \ge \operatorname{mean}(s)$ (the whole interface at its
  threshold is the largest transmissible mean traction); the solver rejects
  such targets up front rather than diverging.
\end{itemize}

\begin{algorithm}[t]
\caption{Force control: outer Newton/Broyden on $\bdelta_t$
  (\code{solve\_tangential})}
\label{alg:force}
\KwIn{target mean traction $\bar\bq$ with
  $\abs{\bar\bq} < \operatorname{mean}(s)$; optional carried stiffness
  $\bK$, warm start, history}
\lIf{no valid $\bK$ carried}{initialise $\bK$ from two
  $\varepsilon$-probe inner solves (differential at the operating point if
  history is present)}
$\bdelta \leftarrow \bK^{-1} \bar\bq$ (or the carried/supplied initial
shift)\;
\For{$\mathrm{outer} = 1, \dots, 40$}{
  $\bq \leftarrow$ inner solve (Algorithm~\ref{alg:pcgtan}) at $\bdelta$,
  warm-started\;
  $\bF \leftarrow \operatorname{mean}(\bq) - \bar\bq$; keep best-$\abs\bF$
  iterate\;
  stop if $\abs\bF \le$ tol, or if $\bF$ stopped responding to $\bdelta$
  (inner floor reached)\;
  Broyden: $\bK \mathrel{+}= \bigl((\Delta\bF - \bK\,\Delta\bdelta)\,
    \Delta\bdelta^{\tr}\bigr)/\abs{\Delta\bdelta}^2$,
    \emph{rejected} if it degenerates $\det\bK$\;
  $\bdelta \mathrel{-}= \bK^{-1}\bF$\;
}
terminal correction: distribute the residual load over interior points,
re-clamp (2--3 passes $\to$ exact load)\;
\Return best iterate, its $\bdelta$, updated $\bK$ (carried to the next
step)\;
\end{algorithm}

% ---------------------------------------------------------
\section{Load Stepping: the Friction Driver}
\label{sec:driver}

\subsection{State and history}
\label{sec:warm}

History dependence (\S\ref{sec:toy1}) means the solver must carry state
between steps. \code{FrictionDriver} owns:

\begin{center}
\begin{tabular}{lll}
\toprule
state & meaning & role \\
\midrule
$p$ & pressures at last accepted step & normal warm start; thresholds \\
$\bq$ & total tangential tractions & the history in \eqref{eq:qp}; warm start \\
$\bu_t = \bC\bq$ & tangential displacements & avoids one rebuild matvec per step \\
$\bdelta_t$ & accumulated rigid shift & converts total targets to increments \\
$\bw_{\mathrm{acc}}$ & accumulated slip & output/diagnostics (e.g.\ fretting maps) \\
$\bK$ & carried force-control stiffness & skips the $\varepsilon$-probes \\
\bottomrule
\end{tabular}
\end{center}

The incremental QP \eqref{eq:qp} is fed to the solver in \emph{total}
variables through one device: since
$\nabla J = \bC\bq - \bC\bq^n - \Delta\bdelta_t$, passing the
\textbf{history offset} $\bu_{\mathrm{hist}} = -\bC\bq^n = -\bu_t^n$ into
the residual, $\bg_i = (\bC\bq)_i + \bu_{\mathrm{hist},i} -
\Delta\bdelta_t$, makes Algorithm~\ref{alg:pcgtan} solve the incremental
problem while iterating directly on the total traction $\bq$ (so the disk
constraints, which act on \emph{totals}, are natural). User-facing targets
are totals; the driver converts to increments internally.

Warm starts do most of the practical work in a load program: the previous
step's $\bq$ seeds the CG, the previous shift seeds $\bdelta$, and the
carried $\bK$ skips the probe solves --- measured on a repeated force-solve:
$1475 \to 14$ inner iterations. A warm-started, nearly converged step would
pay the full stagnation window just to certify itself; the driver floors the
error normalisation ($g_{\mathrm{floor}}$, from the history scale) so such
steps exit immediately. Note the unit discipline: under force control the
target is a \emph{traction} and must not enter the
displacement-residual floor (unit mixing at non-unit $\Estar$).

\subsection{The transactional step}
\label{sec:transactional}

\code{step()} computes everything on local candidates and commits the
persistent state \emph{only if the whole step converged} (normal solve,
threshold loop, tangential solve). A non-converged step returns
\code{converged = false} and leaves the driver exactly as before the call
--- the caller may retry with a smaller increment, a looser tolerance, or
\code{reset()}. This mirrors transactional semantics in databases and
spares the user the subtle corruption of half-updated friction history
(which, being path-dependent, could not be repaired afterwards).

\begin{algorithm}[t]
\caption{One driver step (\code{FrictionDriver::step})}
\label{alg:driver}
\KwIn{targets: mean pressure $\bar p$ and/or total $\bar\bq$ \emph{or}
  total $\bdelta_t$; step duration $\Delta t$; optional field $T$}
\If{$\bar p$ given}{
  $p \leftarrow$ Polonsky--Keer PCG (warm start: previous $p$);
  abort (unchanged state) if not converged\;
}
\If{tangential target given}{
  target increment $\leftarrow$ total target $-$ carried total;\quad
  $\bu_{\mathrm{hist}} \leftarrow -\bu_t^n$\;
  \Repeat{$s$ change $<$ \code{threshold\_rtol}
    \emph{(single pass for rate-independent laws)}}{
    $v_i \leftarrow \abs{\Delta\bw_i}/\Delta t$ (pass 0: previous step's
    slip)\;
    $s \leftarrow \tfrac12 (s + \tau_c(p, v, T))$ (pass 0: direct)\;
    $\bq \leftarrow$ \code{solve\_tangential} (Algorithms
    \ref{alg:pcgtan}/\ref{alg:force}; warm-chained across passes)\;
    $\Delta\bw_i \leftarrow \Delta\bdelta_t - \bigl(\bC(\bq -
    \bq^n)\bigr)_i$\;
  }
  $D \leftarrow h^2 \sum_i \bq_i \cdot \Delta\bw_i$
  \Comment{report; tests assert $D \ge 0$}
}
\textbf{commit} $p, \bq, \bu_t, \bdelta_t, \bw_{\mathrm{acc}}, \bK$
\emph{only if everything converged}\;
\end{algorithm}

% ---------------------------------------------------------
\section{The Bipotential Uzawa Reference Solver}
\label{sec:uzawa}

\subsection{Why a second, completely different solver}

The staggered path of \S\ref{sec:solver} is fast but rests on a stack of
choices (staggering, two-metric projection, outer Newton) that could, in
principle, all be consistently wrong together. The project therefore
maintains an independent reference solver built on entirely different
mathematics --- the \emph{bipotential} method of de~Saxc\'e \&
Feng~\cite{DeSaxceFeng1998} --- and gates releases on the two families
agreeing. A second motivation is strategic: Coulomb friction is a
\emph{non-associated} law (the slip direction has no component ``opening''
the contact, unlike what a normal-dissipation postulate on the cone would
predict), so the coupled normal--tangential problem does not derive from a
single convex minimisation. The bipotential framework was designed
precisely for such laws, treats normal contact and friction
\emph{simultaneously} through one implicit rule, is derivative-free, and is
indifferent to operator symmetry --- making it the designated strategy for
the future coupled (dissimilar-material) case, where staggering loses its
justification.

\subsection{The algorithm: predictor--corrector with a cone projection}

Per sweep, for every point, the traction triple $(q_x, q_y, p)_i$ is
updated by a step of an Uzawa iteration
(eqs.~(102)--(105) of \cite{DeSaxceFeng1998}, transplanted onto the BEM
operators in \code{bipotential.cpp}): an \emph{augmented predictor} that
mixes the current traction with the current kinematic violation, followed
by an exact projection onto the friction cone. In this codebase's
conventions (the sign map is derived in \code{bipotential.hpp}: the paper's
relative velocity is $-\bw$ and its separation velocity is the gap
$g = u_z + g_0 - \alpha \ge 0$),
\begin{equation}
  \btau_{t} = \bq_i + \rho\, \bw_i,
  \qquad
  \tau_n = p_i - \rho \left( g_i + \mu \abs{\bw_i} \right),
  \label{eq:predictor}
\end{equation}
followed by the analytic orthogonal projection of
$(\btau_t, \tau_n)$ onto $K_\mu$, which has three mutually exclusive cases
with transparent physical meaning:
\begin{equation}
  (\bq_i, p_i) \leftarrow
  \begin{cases}
    (\bzero,\, 0), & \mu\abs{\btau_t} < -\tau_n
      \quad\text{(inside the polar cone: separation)},\\[4pt]
    (\btau_t,\, \tau_n), & \abs{\btau_t} \le \mu\, \tau_n
      \quad\text{(inside the cone: stick)},\\[4pt]
    \left( k\mu\, \hat\btau_t,\ k \right),\;
    k = \dfrac{\tau_n + \mu \abs{\btau_t}}{1 + \mu^2},
      & \text{otherwise}
      \quad\text{(radial return to the cone: slip)}.
  \end{cases}
  \label{eq:cone}
\end{equation}
The $\mu\abs{\bw}$ term inside the \emph{normal} predictor
\eqref{eq:predictor} is the bipotential signature --- it is what encodes
the non-associated coupling (the ``implicit subnormality rule'') and
distinguishes the method from solving an associated cone law. The same
projection formula is used by Tamaas
(\code{Kato::enforcePressureCoulomb}), which was one of the cross-checks
that pinned the constants. Tresca and generic laws replace the cone by a
\emph{cylinder}: $p \ge 0$ and $\abs{\bq_i} \le s_i$ project independently,
with $s$ frozen per iteration.

The step size $\rho$ must satisfy an inverse-operator-norm bound for
convergence; it is set as $\rho = \rho_{\mathrm{scale}} / \lambda_{\max}$
with $\lambda_{\max}$ estimated by 20 power iterations on the stacked
operators $\operatorname{diag}(\bC, \bS)$ (both symbols peak at the lowest
wavenumber, so power iteration converges in a few sweeps).

\subsection{Behaviour and the cross-check gate}

The reference is a first-order method: expect thousands of sweeps
(each sweep is one normal plus one tangential matvec). Measured at
$N_s = 32$, Hertz + Coulomb, partial slip: 790 sweeps to a relative update
of $10^{-9}$, agreeing with the production staggered path to relative
$L_2$ of $5.8\times10^{-8}$ in $p$ and $4.3\times10^{-6}$ in $\bq$, and
passing an independent direct KKT self-check at the $10^{-10}$ level.
That two solver families with disjoint mathematics and disjoint failure
modes agree to these levels is the strongest correctness statement in the
friction test suite. (On an all-slip uniform-threshold Tresca case the
cross-check gap of $7.3\times10^{-3}$ is the \emph{production} side's
documented metric floor, not the Uzawa's --- identified by checking which
side moves when tolerances tighten.) The reference is $\sim 10\times$ the
production wall time at $N_s = 32$ and is not intended for production runs.

% ---------------------------------------------------------
\section{Analytical Benchmarks}
\label{sec:benchmarks}

Frictional contact is blessed with a family of exact solutions that probe
different parts of the implementation. Each is a permanent gate in
\code{test\_friction} / \code{test\_driver}.

\subsection{Cattaneo--Mindlin: the partial-slip sphere}
\label{sec:cm}

The flagship problem~\cite{Cattaneo1938,Mindlin1949,Johnson1985}: a sphere
pressed by normal force $P$ (Hertz contact of radius $a$, peak pressure
$p_0$), then pushed tangentially by $Q < \mu P$. Slip starts at the contact
edge (where the pressure --- hence the threshold $\mu p$ --- vanishes) and
an annulus $c \le r \le a$ slips while the core $r < c$ sticks, with
\begin{equation}
  \boxed{\;
  \frac{c}{a} = \left( 1 - \frac{Q}{\mu P} \right)^{1/3},
  \;}
  \qquad
  q(r) = \mu\, p_{\mathrm{Hertz}}(r; a)
       - \mu\, \frac{c}{a}\, p_{\mathrm{Hertz}}\!\left(r; c\right)
       \ \ (r < c),
\end{equation}
i.e.\ the shear traction is the difference of two scaled Hertz pressure
distributions --- the simplest instance of the superposition idea of
\S\ref{sec:cj}. As $Q \to \mu P$, $c \to 0$: gross sliding. Measured
($N_s = 128$, $\nu = 0$, $Q = 0.5\,\mu P$): stick radius $c/a = 0.7972$ vs
the analytic $0.7937$, traction field relative $L_2$ error $0.0074$, and
$q_y$ \emph{exactly} zero (at $\nu = 0$ the $xy$ coupling kernel vanishes
identically, so any non-zero $q_y$ would flag a parity bug).

\subsection{Mindlin's full-stick tangential stiffness}

Before any slip, a circular contact of radius $a$ loaded tangentially
responds linearly~\cite{Mindlin1949}:
\begin{equation}
  Q = \frac{8 G a}{2-\nu}\, \delta_x ,
\end{equation}
the tangential analogue of the normal contact stiffness $2\Estar a$ (their
ratio is the compliance ratio already met in
Example~\ref{ex:selfterm}). Measured at $N_s = 128$, $a = 0.25L$,
$\nu = 0.3$: stiffness ratio $0.9982$ ($0.18\%$ error, consistent with the
grid resolution of the contact edge).

\subsection{Ciavarella--J\"ager superposition and rough surfaces}
\label{sec:cj}

Ciavarella~\cite{Ciavarella1998} and J\"ager~\cite{Jager1998} discovered
independently that for \emph{monotonic} tangential loading of a Coulomb
interface, the shear traction is the difference of two \emph{normal}
solutions of the same geometry:
\begin{equation}
  \bq = \mu \left[ p(F_z) - p^{*}(F^{*}) \right],
  \qquad
  F^{*} = F_z - \frac{F_x}{\mu},
  \label{eq:cj}
\end{equation}
where $p(F_z)$ is the pressure at the actual normal load and $p^*$ is the
pressure the \emph{same} normal problem would produce at the reduced load
$F^*$; the stick zone is exactly the contact zone of the reduced problem
(so $A_{\mathrm{stick}}/A_0 = 1 - F_x/(\mu F_z)$ under the appropriate
load-area law~\cite{PaggiPohrtPopov2014}). The result holds when the
tangential and normal influence operators are proportional and uncoupled
--- on the discrete grid, exactly at $\nu = 0$, where by
\S\ref{sec:symbol} $\bC = \frac{1}{1-\nu}\,\operatorname{diag}(\bS,\bS)
\big|_{\nu=0} = \operatorname{diag}(\bS, \bS)$ up to the constant
absorbed in \eqref{eq:prefactor}, and $\bC_{xy} \equiv 0$.

This yields validation of \emph{discrete-exact} type, the most valuable
kind: the friction solver's output is compared against two runs of the
long-validated \emph{normal} solver on the same grid --- no discretisation
error enters the comparison, only solver error. Measured: $q_x$ relative
$L_2 = 3.2\times10^{-6}$ on a rough surface at $N_s = 64$ (gate $10^{-4}$),
with the stick set matching the reduced contact set point-for-point
(206 = 206). The same mechanism validates the \emph{driver}: two-step
loading vs one-step vs \eqref{eq:cj} agree to $10^{-5}$--$4\times10^{-5}$
(monotonic path independence), and a Mindlin-type unloading sequence
matches its own discrete superposition (load $\bq_1$, reverse to $\bq_2$:
the unloading solution is $\bq_1 - 2\mu[p - p^{**}]$-type with counter-slip
from the edge) to $8.8\times10^{-5}$ with $D > 0$ at every step.

\subsection{Limits worth memorising}

\begin{itemize}
\item $\tau_c \to \infty$ (or $\mu \to \infty$): full stick everywhere ---
  the tangential flat-punch analogue; checks the pure elastic path.
\item $Q \to \mu P$: gross sliding onset; the force-control guard of
  \S\ref{sec:force} triggers exactly at the mean-threshold bound.
\item $\mu = 0$: zero shear, and the normal solution must be bit-identical
  to the frictionless solver (regression anchor).
\item $\nu = 0$: tangential problem $\equiv$ two copies of the normal
  problem; $q_y \equiv 0$ under $x$-loading.
\end{itemize}

% ---------------------------------------------------------
\section{Validated Numbers and Practical Notes}
\label{sec:numbers}

\subsection{Friction-specific validated numbers}

(The project-wide table lives in \code{CLAUDE.md}; this subset is the
friction feature's permanent record. All measured 2026-07 on the 20-core
reference workstation.)

\begin{center}
\small
\begin{tabular}{ll}
\toprule
quantity & value \\
\midrule
Cerruti closed forms vs $64^2$ GL quadrature (6 offsets, incl.\ self) & rel.\ err $< 10^{-9}$ \\
tangential FFT matvec vs dense ($N_s=8/16$) & $2.9\times10^{-16}$ rel.\ $L_2$ \\
tangential H$^2$ matvec vs dense ($N_s=32$, $q{=}4$ / $q{=}6$) & $1.67\times10^{-4}$ / $6.97\times10^{-6}$ \\
KKT, displacement control ($N_s=32$, 28 slip pts, tol $10^{-5}$) & 21 iterations; stick $\abs{\bw}$ $7.7\times10^{-7}$ \\
force control ($N_s=32$, $40\%$ slip, outer Broyden) & load met to $3\times10^{-16}$ rel. \\
preconditioner A/B ($N_s=64$ partial slip) & $271 \to 126$ iterations ($2.15\times$) \\
Mindlin full-stick stiffness $Q = 8Ga\delta/(2-\nu)$ ($N_s=128$) & ratio $0.9982$ \\
Cattaneo--Mindlin $\nu=0$ ($N_s=128$, $Q=0.5\mu P$) & $c/a$ $0.7972$ vs $0.7937$; $q_y \equiv 0$ \\
Ciavarella--J\"ager, rough, $\nu=0$ ($N_s=64$) & $q_x$ rel.\ $L_2$ $3.2\times10^{-6}$; stick set exact \\
force-control $\bK$ carry-over (repeat solve) & $1475 \to 14$ inner iterations \\
driver two-step C--J ($N_s=64$, rough, Coulomb) & rel.\ $4.2\times10^{-5}$; $D > 0$ \\
monotonic path independence (2-inc / 1-inc / C--J) & $1.4$/$4.0$/$3.1 \times10^{-5}$ \\
Mindlin unloading superposition ($\nu=0$, counter-slip) & rel.\ $8.8\times10^{-5}$; $q_y \equiv 0$; $D>0$ \\
velocity-dependent threshold loop (rate-weakening) & 7 damped passes; $D > 0$ \\
bipotential Uzawa vs production ($N_s=32$, Coulomb) & 790 sweeps; $p$: $5.8\times10^{-8}$, $\bq$: $4.3\times10^{-6}$ \\
bipotential KKT self-check & $\sim 10^{-10}$ \\
\bottomrule
\end{tabular}
\end{center}

\subsection{Tolerances and floors (what to expect in practice)}

\begin{itemize}
\item The displacement-control KKT metric floor is
  $\sim 10^{-5}$--$10^{-4}$ relative, growing with the slip fraction (slip
  directions rotate; clamps perturb conjugacy). The default
  \code{tol\_tangential} of $10^{-5}$ is calibrated to it.
\item Under force control the \emph{load} is met to machine precision
  (terminal correction), while the pointwise KKT residual floors at
  $10^{-6}$--$10^{-4}$ (with OpenMP-reduction jitter); the validation gate
  is $10^{-2}$ with the sharp discrete-exact validation delegated to the
  $\nu = 0$ Ciavarella--J\"ager tests.
\item Dissipation is reported per step; tiny negative values at the solver
  floor ($\abs{D} \lesssim$ floor $\times$ load scale) would indicate a
  genuine bug, not noise --- the tests treat $D \ge 0$ strictly.
\end{itemize}

\subsection{Python quick reference}

\begin{verbatim}
import aspher as hc, numpy as np
model = hc.CoulombFriction(mu=0.3)        # TrescaFriction(tau_c=...), UserFriction(fn)
fs = hc.FrictionSolver(grid_size=256, E_star=1.0, nu=0.3, model=model)
fs.set_gap(-surface)                      # gap = -height (rigid flat)
fs.step(p_bar=0.05)                       # normal preload
r = fs.step(q_bar=(0.01, 0.0), dt=1.0)    # force-controlled tangential step
#   fs.step(delta_t=(1e-4, 0.0), dt=1.0)  # displacement-controlled variant
print(r.n_stick, r.n_slip, r.q_mean, r.dissipation)
\end{verbatim}

State is carried across steps; a non-converged step leaves it unchanged
(\S\ref{sec:transactional}). \code{hc.solve\_bipotential(...)} exposes the
reference solver of \S\ref{sec:uzawa} for cross-checks.

% ---------------------------------------------------------
\section*{Acknowledgement of sources}

The element-integrated kernels follow Pohrt \& Li~\cite{PohrtLi2014} and
Dydo \& Busby~\cite{DydoBusby1995}; the incremental variational framing is
classical (see Johnson~\cite{Johnson1985} for the mechanics and
Nocedal--Wright~\cite{NocWri2006} for the optimisation); the projected-CG
skeleton extends Polonsky \& Keer~\cite{PolKeer1999}; the two-metric
projection idea is Bertsekas'~\cite{Bertsekas1982}; the reference solver is
de~Saxc\'e \& Feng's~\cite{DeSaxceFeng1998}; the benchmark solutions are
due to Cattaneo~\cite{Cattaneo1938}, Mindlin~\cite{Mindlin1949},
Ciavarella~\cite{Ciavarella1998}, J\"ager~\cite{Jager1998}, and, for the
rough-surface form, Paggi, Pohrt \& Popov~\cite{PaggiPohrtPopov2014}.

\bibliographystyle{unsrt}
\bibliography{references}

\end{document}
